\textbf{۳۵۵}

گبی گازی به مافین زد و سپس آن را دوباره روی بشقاب گذاشت. نمی‌توانستم بفهمم چه کسی او را به خانه دعوت کرده بود. موریک دور و بر نبود و او با هیچ کس دیگر اینجا دوستی نداشت.

گفتم: «باید کار کنم. اما از آشنایی با تو خوشحال شدم.»

تصور می‌کردم که خودش راه خروج از خانه را پیدا می‌کند. اما او باید اشتباهی کرده بوده باشد. مرموز بعداً او را در حالی کشف کرد که روی توالت او نشسته بود. اولین باری که او را دیدم، داشت دستشویی می‌کرد. هر دو چنان نارسیسیست بودند که فکر می‌کردم مانند دو انتهای مثبت آهن‌ربا همدیگر را دفع کنند. اما به جای آن، با هم رابطه داشتند.

پرسیدم: «تو کی هستی؟» او هفته بعد را در خانه گذراند، با مرموز خوابید و با کردن کورتی شکایت می‌کرد چون لباس‌هایش را بدون اجازه قرض می‌گرفت.

گبی از دوستان موریک بود، یکی از بسیاری از جوانان \lr{PUA} که هر هفته غیرمجاز به خانه ما می‌آمدند و در اتاق نشیمن ظاهر می‌شدند. او نگرش یک ملکه زیبایی را داشت اما بدنش مانند یک کیسه گوجه‌فرنگی بود. قدمی به عقب برداشتم و شروع به بستن در کردم.

«هی»، او گفت و توالت را شست. «این خانه قشنگیه. تو چیکار می‌کنی؟»

آن کلمات بلافاصله همه چیز را خراب کرد. وقتی در لس‌آنجلس دختران را سارج می‌کنید، راداری برای زنانی که استفاده‌گر هستند، ایجاد می‌کنید. آنها که بی‌تاکتیک‌تر هستند، در دقایق ابتدایی گفتگو می‌پرسند چه ماشینی دارید، یا چه کار می‌کنید، یا با کدام یک از افراد مشهور حاضر در اتاق دوست هستید تا رتبه اجتماعی شما و میزان مفید بودنتان را تعیین کنند. آنهایی که بی‌ تکلف نیستند، نیازی به پرسش ندارند: به ساعت شما نگاه می‌کنند؛ می‌بینند مردم چطور به شما پاسخ می‌دهند وقتی صحبت می‌کنید، به دنبال نشانه‌هایی از عدم اعتماد به نفس در گفتارتان گوش می‌دهند. اینها نشانه‌هایی هستند که \lr{PUAs} به آن‌ها «زیر ارتباط» می‌گویند.

گبی به دسته بی‌تکلف‌تر این گونه تعلق داشت. 

در حال شستن دست‌هایش، او در کابینت دارویی را باز کرد و محتویات آن را بررسی کرد. سپس وارد اتاق من شد و به جستجویش ادامه داد. «نویسنده‌ای؟» پرسید. «باید درباره من بنویسی. داستان خیلی جالبی دارم. می‌خوام بازیگر بشم. می‌دونی بعضی از مردم ذاتاً برای شهرت به دنیا میان.» یک جفت عینک آفتابی روی میز لباس من برداشت و آن را به چشم زد. «خب، من یکی از اون‌ها هستم. نه اینکه بخوام بگم خاص هستم یا چیزی. فقط چیزی هست که از سن کم می‌دانی چون مردم با تو متفاوت برخورد می‌کنند.»

یک مرد ثروتمند نیازی ندارد بگوید ثروتمند است.

همان‌طور که او مشغول صحبت بود، یک مافین از بشقابی روی میز من برداشت. امروز روز مافین بود. کورتی در خانه می‌دوید و به همه بشقاب‌هایی پر از مافین می‌داد به طوری که بیشتر از آن می‌توانستند بخورند.

گبی در حالی که دستانش را می‌شست، کابینت دارویی را باز کرد و محتویاتش را بررسی کرد. سپس وارد اتاق من شد و به گشت‌وگذار خود ادامه داد. «تو نویسنده‌ای؟» پرسید. «باید درباره من بنویسی. داستان خیلی جالبی دارم. می‌خوام بازیگر بشم. بعضی‌ها ذاتاً برای شهرت زاده می‌شن.» یک جفت عینک آفتابی از روی بالا کشوی من برداشت و به چشم زد. «خب، من یکی از اونام. نه اینکه بخوام بگم خاص‌هستم یا چیزی. فقط چیزیه که از یک سن کم می‌دونی، چون مردم با تو متفاوت رفتار می‌کنن.»

مرد ثروتمند نیازی نداره به تو بگه که ثروتمنده. در حالی که مدام حرف می‌زد، یک مافین از بشقاب روی میز من برداشت. امروز روز مافین بود. کورتی در خانه می‌دوید و به همه بشقاب‌هایی پر از مافین بیشتری از آنچه که بتوانند بخورند، می‌داد.

بعدازظهری، وقتی کورتی در آشپزخانه با دو قاشق کره بادام‌زمینی توی ظرف کنار مافین‌ها ایستاده بود، از گبی پرسید: «نمی‌خوای خونه بری؟»

گبی با حالتی عجیب به او نگاه کرد و گفت: «خونه؟ من همین‌جا زندگی می‌کنم.»

این موضوع برای کورتی، برای من و برای مرموز نیز تازه بود. خانه چنین آدم‌هایی را به خود جذب می‌کرد؛ در نهایت همه را بیرون می‌راند.

تویلا قربانی بعدی پروژه هالیوود بود. اولین بار زمانی در خانه ظاهر شد که یک رقاصه که مرموز سال‌ها پیش با او معاشرت کرده بود، در حالی‌که دچار افسردگی شدیدی شده بود جلوی در خانه آمد. با داشتن تجربه‌ای در این زمینه، مرموز پیشنهاد کرد که شبی به او مشاوره بدهد وقتی گبی بیرون با دوستانش بود. 

اما آن رقاصه مست سر رسید و همراهش تویلا را هم آورد.

تویلا جایزه‌ای نبود. او یک هالیوودی راک و رولر سی و چهار ساله با پوستی چروکیده، بدنی سفت مثل صورتش، موهایی مشکی به شکل لانه پرنده و قلبی از طلا بود. او مرا به یاد یک پونتیاک فیرو می‌انداخت، یک مدل قدیمی اسپرت که امکان خرابی‌اش هر لحظه وجود داشت.

وقتی مرموز و تویلا شروع به ناز و عشوه کردند، دوست مست و افسرده‌شان به گریه افتاد. نیم‌ساعت در چاله بالشتی گریه کرد تا وقتی تویلا و مرموز پا به اتاق او گذاشتند. گبی آن شب به خانه برگشت و بدون یک کلمه اعتراض، با آن دو به تخت افتاد و بلافاصله خوابش برد. گبی و مرموز عاشق یکدیگر نبودند؛ آنها فقط سرپناه یکدیگر را می‌خواستند.

آن صبح و صبح بعد، تویلا برای همه در خانه پنکیک پخت. چون به نظر نمی‌آمد قصد ترک خانه را داشته باشد، مرموز او را به عنوان دستیار شخصی خودش با حقوق چهارصد دلار در هفته استخدام کرد.